% =============================================================================
\chapter{Finite Elements, Degrees of Freedom, and Basis Functions}
\label{ch:elements}
% =============================================================================

This chapter introduces the local building blocks of finite element spaces. It
starts from the Ciarlet definition of an element, then develops degrees of
freedom, basis construction, interelement continuity, and the principal element
families used in practice.

\section{The Ciarlet definition of a finite element}
\label{sec:elements:ciarlet}

% Cell, local approximation space, and degrees of freedom

\section{Degrees of freedom and unisolvence}
\label{sec:elements:unisolvence}

% Interpolation conditions, nodal variables, moments, and dual bases

\section{Scalar basis functions}
\label{sec:elements:scalar_basis}

% Nodal, modal, and hierarchical bases; partition of unity and locality

\section{Classical element families}
\label{sec:elements:families}

% Lagrange, serendipity, bubble-enriched, spectral, and isogeometric variants

\section{Vector-valued and compatible elements}
\label{sec:elements:compatible}

% Raviart-Thomas, Brezzi-Douglas-Marini, Nedelec, and exact-sequence ideas

\section{Transformations and interelement continuity}
\label{sec:elements:transforms}

% Scalar pullbacks, Piola transforms, orientation, and conformity requirements

\section{Global finite element spaces}
\label{sec:elements:global_spaces}

% Assembly of local shape functions into conforming and nonconforming spaces
